对协方差自适应后端而言，这种失配直接转化为求解困难。以 CMA-ES 为例：当同一子问题内同时包含极平坦与极陡峭的方向时，优化器陷入步长两难——足以沿平坦谷底推进的步长对陡峭方向过大，而对陡峭方向不致发散的步长又使平坦方向推进极慢。协方差自适应机制能逐步缓解这种失配，但需要充足的采样与学习。在有限预算下，分组所决定的局部条件性可实质性地增加后端求解难度；相比之下，现有分组研究对这一优化器实际面对的局部条件性着墨甚少。